\section{Reconstructing the partitions on the original manifold}
\label{sec:reconstruction}

We prove \Cref{topology:partitions}, the partition statement used in
Section~\ref{sec:width-map}. Keep the original metric $g$, its distance
$d_0$, the centers, and the ball pairs fixed there. The finite-stage
separating sets and proper collapses are supplied by \Cref{countable:construction}.
Only finite truncations of the analytic product ends enter the
reconstruction. The choice of homeomorphisms below makes every change
for a center label lie in its original ball of radius $2D+2$; local
finiteness of those balls then gives the limit on one copy of $M$.

\begin{proof}[Proof of \Cref{topology:partitions}]
Choose positive numbers $\eps_0,\eps_1,\ldots$ with
$\sum_{r\ge0}\eps_r<1$, and use $\eps_j$ as the error of
\Cref{component:small-displacement} for the split of label $j$, the transition
from analytic stage $j$ to $j+1$ in \Cref{countable:construction}.  In this
proof only, $X_r$ denotes the reconstruction after the radial label $0$ and
the center labels $1,\ldots,r$ have been split; it is $X_{r+1}$ in the
indexing of \Cref{countable:construction,component:reconstruction}.  Thus the
passage from $X_r$ to $X_{r+1}$ splits the center label $i=r+1$, with
error $\eps_i=\eps_{r+1}$.
For a finite reconstruction write
\[
 q_r:X_r\longrightarrow M,\qquad
 H_r:X_r\longrightarrow M,\qquad
 Q_r=q_r\circ H_r^{-1}.
\]
The map $q_r$ is the proper collapse of
\cref{component:reconstruction}, whereas $H_r$ is a homeomorphism to a
fixed copy of the original manifold.  We maintain
\begin{equation}\label{topology:finite-error}
 d_0(Q_r(x),x)\le E_r,
 \qquad E_r\le\sum_{j=0}^{r}\eps_j<1.
\end{equation}
Separating sets and their two closed sides will also be viewed on this fixed copy
using $H_r$.

For the radial stage, the cut has a locally finite family of compact
components.  Choose pairwise disjoint, locally finite collars of these
components on which the collapse varies by less than $\eps_0$ in each
normal segment.  Compactness of each component and continuity of the
collapse permit the collars to be this small, independently from one
component to another.  Compress the finite inserted cylinders into
these collars.  Since their supports are disjoint, the displacement
bound is $\eps_0$, rather than a sum over components.  This constructs
$H_0$ and proves \eqref{topology:finite-error} initially.  The finite
length of an inserted cylinder may be arbitrarily large; its length
does not impose a lower bound on the width of the collar into which it
is compressed.

Fix $r\ge0$ and put $i=r+1$. First raise any old finite
truncation levels needed for the new compatible defining functions.  Denote the
resulting reconstruction by $\widehat X_r$.  The truncation level
reparameterization
\[
 \Phi_r:\widehat X_r\longrightarrow X_r
\]
preserves collapse and every old separating set and signed side.  Indeed, on an
old product end the collapse is independent of the normal coordinate,
and an increasing change of that coordinate preserves its sign and its
zero section.  Compatible products of these changes agree at corners.
Thus we may set
\begin{equation}\label{topology:horizon}
 \widehat q_r=q_r\circ \Phi_r,\qquad
 \widehat H_r=H_r\circ \Phi_r,
 \qquad
 \widehat q_r\circ\widehat H_r^{-1}=Q_r.
\end{equation}
In particular, increasing a truncation level contributes no displacement on the
fixed copy.

The new separating set lies over $B_0(x_i,2D)$.  Choose its collar strip $S_i$ in
$\widehat X_r$ so that
\begin{equation}\label{topology:strip}
 S_i\subset\widehat q_r^{-1}\bigl(B_0(x_i,2D+1)\bigr),
\end{equation}
and so that collapse oscillation along every segment used by the
compression is less than the error assigned to this stage.  These
choices are possible because the relevant separating set components in a finite
reconstruction are compact and locally finite.  More specifically,
their union lies in the inverse image of the compact ball
$\overline B_0(x_i,2D)$ under the proper finite collapse; hence only
finitely many components are involved in this label stage.

Choose these strips using the collars in the proof of
\Cref{component:reconstruction}. That construction applies after the
truncation levels have been raised beyond the product thresholds: the new
defining function is regular on each old end base and independent of its
old signed normal coordinates. Its collar flow and the resulting
compression preserve every old separating set and signed side and their
intersections, including the radial ones. Decrease the componentwise strip widths to satisfy
\eqref{topology:strip} and the prescribed collapse oscillation, as
permitted by \Cref{component:small-displacement}.
Let
\[
 c_r:X_{r+1}\longrightarrow\widehat X_r
 \quad\hbox{and}\quad
 h_r:X_{r+1}\longrightarrow\widehat X_r
\]
be, respectively, the new cylinder collapse and the homeomorphism that
compresses those cylinders into the chosen strips.  They agree off
the inverse image of $S_i$. Figure~\ref{fig:topology-compression} shows
the two maps schematically in a normal coordinate: $c_r$ collapses the
inserted interval, while $h_r$ compresses it into the chosen old collar.
Define
\begin{equation}\label{topology:recursion}
 q_{r+1}=\widehat q_r\circ c_r,
 \qquad H_{r+1}=\widehat H_r\circ h_r,
 \qquad
 Q_{r+1}=\widehat q_r\circ c_r\circ h_r^{-1}
                         \circ\widehat H_r^{-1}.
\end{equation}
The oscillation choice gives
$d_0(Q_{r+1}(x),Q_r(x))<\eps_{r+1}$ wherever they differ, proving
\eqref{topology:finite-error} at the next stage.  Their difference is
supported in $\widehat H_r(S_i)$.  If
$x=\widehat H_r(z)$ belongs to this set, then
\eqref{topology:horizon}, \eqref{topology:strip}, and
\eqref{topology:finite-error} imply
\begin{equation}\label{topology:fixed-support}
 d_0(x,x_i)
 \le d_0(x,Q_r(x))+d_0(\widehat q_r(z),x_i)
 <1+(2D+1)=2D+2.
\end{equation}
This support radius is independent of the stage.

\begin{figure}[!b]
\centering
\begin{tikzpicture}[x=1.25cm,y=1.05cm,>=Latex,font=\small]
 \fill[blue!6] (-3,-1) rectangle (3,1);
 \fill[red!8] (-1,-2.25) rectangle (1,2.35);
 \draw[->] (-3.2,0)--(3.5,0) node[right] {$s$};
 \draw[->] (0,-2.3)--(0,2.65) node[above] {$t$};
 \draw[densely dashed,gray] (-3,1)--(3,1);
 \draw[densely dashed,gray] (-3,-1)--(3,-1);
 \draw[very thick,purple!75!black]
  (-3,-2)--(-1,0)--(1,0)--(3,2);
 \draw[thick,orange!85!black]
  (-3,-2)--(-2,-1)--(-1,-.35)--(0,0)--(1,.35)--(2,1)--(3,2);
 \draw (-1,.07)--(-1,-.07) node[below=4pt,fill=white,inner sep=1pt] {$-L$};
 \draw (1,.07)--(1,-.07) node[below=4pt] {$L$};
 \node[left,fill=white,inner sep=1pt] at (-.08,1) {$\delta$};
 \node[left,fill=white,inner sep=1pt] at (-.08,-1) {$-\delta$};
 \node[align=left,anchor=west,text=purple!75!black] at (3.55,1.25)
  {collapse $c_r$};
 \node[align=left,anchor=west,text=orange!85!black] at (3.55,.55)
  {compression $h_r$};
 \node[align=left,anchor=west] at (3.55,-.55)
  {$c_r=h_r$ outside the strip};
 \node[align=left,anchor=west] at (3.55,-1.35)
  {old strata follow\\the same reparameterization};
 \node[below,align=center] at (0,-2.4)
  {finite inserted cylinder: $-L\le s\le L$};
\end{tikzpicture}
\caption{The two maps used in a finite reconstruction, shown schematically
in local normal coordinates: $s$ on the reconstructed domain $X_{r+1}$
and $t$ on the old target $\widehat X_r$. The target strip $|t|<\delta$
represents the chosen collar $S_i$.
The collapse is constant on the inserted cylinder;
the homeomorphism squeezes it into an arbitrarily narrow old collar.
Where the maps differ, both images lie in the target strip.  Infinite analytic ends
are truncated before this topological operation.}
\label{fig:topology-compression}
\label{component:collar-figure}
\end{figure}

The balls in \eqref{topology:fixed-support} form a locally finite
family.  Consequently, on every relatively compact neighborhood in
the fixed copy, the maps $Q_r$ are exactly unchanged after finitely
many stages.  Their locally stationary limit $q$ is continuous and
satisfies \eqref{topology:displacement}.  It is proper: the inverse
image of a compact set is a closed subset of its compact closed
$1$-neighborhood in $(M,d_0)$.

At each finite stage the signed sides are the closed reconstructed
sides, with their intersection equal to the central separating set.  The
stratum-preserving identifications just constructed preserve every
previously installed side and separating set exactly.  Newly installed label
separating sets are confined to the fixed balls in
\eqref{topology:fixed-support}.  Thus all the partition data are locally
eventually unchanged.  Their final closed sets satisfy
\eqref{topology:sides} and \eqref{topology:forced}, since these are the
finite-stage forced-sign properties of
\cref{countable:construction}, and both collapse and indices are
stationary near each point.  The separating set family is locally finite by the
same support bound.  Finally, a point on $n-1$ distinct final label
separating sets would, on a neighborhood where these data have stabilized, lie
on the same $n-1$ finite-stage separating sets.  This contradicts
\cref{countable:construction}.
\end{proof}
